\section{数据源与相关性测量}

\begin{frame}{四源不是同质搜索引擎：每个源承担不同证据通道}
\scriptsize
\begin{tabularx}{\textwidth}{p{0.13\textwidth} p{0.16\textwidth} p{0.12\textwidth} X X}
\toprule
源 & 通道 / trust & 自然键 & 当前可用算子 & 最适用任务 \\
\midrule
arXiv & research / 0.72 & arXiv ID & 字段布尔、category、submittedDate、start/max & 新研究、攻击/防御概念、趋势发现 \\
NVD & advisory / 0.90 & CVE ID & keyword/exact、CWE、CVSS、CPE、hasKev、发布/修改窗、分页 & 产品漏洞、CVE/CWE/CVSS 结构化确认 \\
OSV & advisory / 0.84 & OSV/GHSA ID & ID、package、ecosystem、purl、version & 开源依赖和模型/agent 工具链供应链 \\
CISA KEV & exploitation / 0.92 & CVE ID & CVE、vendor/product、关键词 & 已知在野利用的确认与优先级 \\
\bottomrule
\end{tabularx}

\vspace{0.6em}
\begin{block}{source-aware 的含义}
同一主题不能要求所有源都提供同一种证据。adaptive coverage 只为 research 和适用的 advisory 建 target；CISA 是 confirmation 通道，不被错误地当作每个主题都必须命中的配额源。
\end{block}
\smallsource{当前映射：\pathcode{engine/gaps.py:28-62}；官方能力：\cite{nvdapi,cisakev,osvapi,arxivapi}}
\end{frame}

\begin{frame}{当前连接器的真实能力与边界}
\small
\begin{columns}[T,onlytextwidth]
\column{0.48\textwidth}
\begin{exampleblock}{与官方语义一致的部分}
\begin{itemize}
  \item NVD 支持 CVE/CWE/CVSS、日期、hasKev、exact 与分页
  \item arXiv 保留主题 query，再注入权威 submittedDate
  \item CISA 可做 CVE 精确确认
  \item 429、5xx、网络错误最多重试 2 次
\end{itemize}
\end{exampleblock}
\column{0.48\textwidth}
\begin{alertblock}{必须写进局限的部分}
\begin{itemize}
  \item OSV 不是全文主题搜索；当前“泛查”遍历固定 15 个 AI package
  \item CISA 的高 trust 只表示 exploitation 证据，不表示主题覆盖广
  \item arXiv 吞吐高，但论文不是漏洞在野证据
  \item NVD 对 LLM 原生行为风险存在 CVE 表达缺口
\end{itemize}
\end{alertblock}
\end{columns}
\smallsource{\pathcode{backend/intel_agent/sources.py:44-535}}
\end{frame}

\begin{frame}{Source tool：模型发起调用前后都有确定性复核}
\centering
\resizebox{0.96\linewidth}{!}{%
\begin{tikzpicture}[
  n/.style={draw,rounded corners,minimum height=0.72cm,align=center,font=\scriptsize,thick},
  a/.style={-{Latex[length=1.7mm]},thick,Navy},node distance=0.45cm
]
\node[n,fill=SoftBlue,minimum width=2.2cm] (cand) {approved\\QueryCandidate};
\node[n,fill=SoftAmber,right=of cand] (identity) {candidate ID\\source/channel};
\node[n,fill=SoftAmber,right=of identity] (time) {权威时间窗\\参数清洗};
\node[n,fill=SoftAmber,right=of time] (budget) {预算/重复\\分页水位};
\node[n,fill=SoftGreen,right=of budget] (http) {注册 client\\HTTP + parser};
\node[n,fill=SoftGreen,right=of http] (out) {去重结果\\QueryCallOutcome};
\draw[a] (cand)--(identity);\draw[a] (identity)--(time);\draw[a] (time)--(budget);\draw[a] (budget)--(http);\draw[a] (http)--(out);
\end{tikzpicture}
}

\vspace{1em}
\begin{columns}[T,onlytextwidth]
\column{0.48\textwidth}
\begin{block}{成功时记录}
raw、in-scope、duplicate、new、relevant-new、has-more、operator signature、checkpoint 证据。
\end{block}
\column{0.48\textwidth}
\begin{block}{失败时记录}
retryable、分页未尽、时间窗不合法等转成 \pathcode{RunGap}，防止把“本轮没查到”误判为“语料已饱和”。
\end{block}
\end{columns}
\smallsource{\pathcode{backend/intel_agent/tools/source_tools.py:90-245}}
\end{frame}

\begin{frame}{三层相关性管线：便宜判断优先，不确定样本升级}
\centering
\resizebox{0.96\linewidth}{!}{%
\begin{tikzpicture}[
  n/.style={draw,rounded corners=2pt,minimum height=0.92cm,minimum width=3.0cm,align=center,font=\small,thick},
  a/.style={-{Latex[length=1.8mm]},thick,Navy}
]
\node[n,fill=SoftBlue] (rule) {Layer 1: 规则\\accept $\ge 0.75$ / reject $\le 0.20$};
\node[n,fill=SoftGreen,right=0.9cm of rule] (emb) {Layer 2: embedding\\accept $\ge 0.62$ / reject $\le 0.42$};
\node[n,fill=SoftAmber,right=0.9cm of emb] (llm) {Layer 3: LLM batch\\relevant=0.85 / irrelevant=0.15};
\draw[a] (rule) -- node[above,font=\scriptsize]{灰区} (emb);
\draw[a] (emb) -- node[above,font=\scriptsize]{灰区} (llm);
\node[below=0.7cm of emb,draw,dashed,rounded corners,fill=Paper,minimum width=8.8cm,align=center,font=\scriptsize] (meta) {输出写入 metadata.relevance + topic\_scores；embedding 以 content hash 缓存};
\draw[a,dashed] (rule.south) |- (meta.west);
\draw[a,dashed] (emb.south) -- (meta.north);
\draw[a,dashed] (llm.south) |- (meta.east);
\end{tikzpicture}
}

\vspace{0.8em}
\begin{alertblock}{配置依赖}
没有 GLM key 或显式禁用时，\pathcode{default_pipeline()} 返回 None；此时 legacy 分析使用原始 \pathcode{relevance_score >= 0.5} 兜底。
\end{alertblock}
\smallsource{\pathcode{backend/intel_agent/relevance.py:31-35,154-205,312-318}；\pathcode{analysis.py:78-82}}
\end{frame}

\begin{frame}{初始 relevance score 的数理语义：它是启发式，不是校准概率}
\small
\begin{block}{source parser 的初始估计}
\[
s_{0}(x)=\min\left(1,\;0.45+\min(0.30,0.10\,n_{topic})+\min(0.20,0.02\,n_{token})\right)
\]
其中 $n_{topic}$ 是匹配主题数，$n_{token}$ 是启发式安全词命中数。CISA 最低取 0.55，
已分类的 NVD item 最低可取 0.68。
\end{block}

\begin{columns}[T,onlytextwidth]
\column{0.48\textwidth}
\begin{exampleblock}{正确用法}
把 $s_0$ 当作低成本先验，交给 rule/embedding/LLM 管线校正，并在 MSA 中测 precision/recall/calibration。
\end{exampleblock}
\column{0.48\textwidth}
\begin{alertblock}{错误用法}
在 relevance pipeline 关闭时直接以 0.5 为 relevant 阈值：只需少量词命中就可能过线，造成 coverage 虚高。
\end{alertblock}
\end{columns}
\smallsource{\pathcode{backend/intel_agent/sources.py:329,373,419,468,537-540}}
\end{frame}

\begin{frame}{测量口径必须按运行配置分层解释}
\scriptsize
\begin{tabularx}{\textwidth}{p{0.23\textwidth} p{0.22\textwidth} X p{0.19\textwidth}}
\toprule
配置状态 & item 判定 & coverage 影响 & 风险 \\
\midrule
有 relevance 标签 & 只有 label=relevant 进入 legacy relevant count & uncertain 不计入 legacy quota & 阈值仍需人工校准 \\
有 metadata score & adaptive 有效证据用 score 连续加权；irrelevant=0 & uncertain 可能贡献部分 evidence & score 未必是校准概率 \\
pipeline 关闭/失败 & fallback: raw score $\ge 0.5$ & 启发式分数可能大量过线 & precision 与 coverage 同时虚高 \\
多主题 item & topic credit 总和最多 1 & 防止一条 item 多格全额计数 & 依赖 topic score 质量 \\
\bottomrule
\end{tabularx}

\vspace{0.25em}
\begin{block}{Measure 阶段的先后关系}
先冻结样本并完成人工双标注、分类 MSA、概率校准与跨源一致性；再比较 bandit、query rewrite 和 stopping。否则控制器只是在优化测量误差。
\end{block}
\end{frame}
